\documentclass[11pt]{article}
\usepackage[margin=1in]{geometry}
\usepackage{amsmath,amssymb,booktabs}
\usepackage{hyperref}
\title{Koopman-Based Linear Embeddings for Control of Automated Agents}
\author{Syed Hafiz Ali}
\date{Research manuscript in preparation}
\begin{document}
\maketitle
\begin{abstract}
We outline an empirical study of finite-dimensional lifted linear models for approximating nonlinear agent dynamics. The accompanying laboratory implements a monomial dictionary and a ridge-regularized EDMD-style operator fit on scalar transition data. The present artifact is a methodological prototype; it does not establish Koopman invariance, control performance, or physical-system validity.
\end{abstract}
\section{Introduction}
Finite-dimensional approximations of the Koopman operator can be constructed from snapshot pairs and observable dictionaries using EDMD \cite{williams2015edmd}. Controlled linear predictors based on lifted representations have also been studied for model predictive control \cite{korda2018mpc}.,The current repository provides a transparent baseline for studying lifted dynamics.
A nonlinear state transition may become easier to analyze after mapping observations into a feature space. A finite dictionary provides a practical approximation, but approximation quality depends on dictionary choice, data coverage, regularization, and the operating region.
\section{Problem formulation}
Given paired observations $(x_t,x_{t+1})$, choose a lifting map $\phi(x)\in\mathbb{R}^d$ and estimate $K$ such that $\phi(x_{t+1})\approx K^\top\phi(x_t)$. The estimate is an empirical finite-dimensional approximation, not proof of an invariant subspace.
\section{Method}
The companion code uses monomials $\phi(x)=(x,x^2,\ldots,x^d)$ and ridge-regularized least squares. It exposes the operator convention, prediction routine, and multi-step lifted rollout.
\section{Experimental protocol}
Compare persistence and fitted linear-in-state baselines against dictionary degrees 1--3. Use chronological train/test splits, held-out state ranges, multiple nonlinear data-generating systems, and fixed seeds. Report one-step and multi-step RMSE, maximum absolute error, sensitivity to regularization, and out-of-range degradation.
\section{Results}
No benchmark results are claimed in this draft. Add complete tables, configurations, and machine-readable outputs only after running the protocol.
\section{Limitations}
A scalar polynomial example is not evidence of a useful embedding for high-dimensional agents. Finite-horizon prediction error is not a control guarantee. Future work must test higher-dimensional systems and close the loop with explicit safety constraints.
\section{Conclusion}
The current repository provides a transparent baseline for studying lifted dynamics. Generalization and control utility remain open empirical questions.
\begin{thebibliography}{9}
\bibitem{williams2015edmd} M. O. Williams, I. G. Kevrekidis, and C. W. Rowley. A Data-Driven Approximation of the Koopman Operator: Extending Dynamic Mode Decomposition. \textit{Journal of Nonlinear Science}, 25:1307--1346, 2015. DOI: \href{https://doi.org/10.1007/s00332-015-9258-5}{10.1007/s00332-015-9258-5}.
\bibitem{korda2018mpc} M. Korda and I. Mezić. Linear Predictors for Nonlinear Dynamical Systems: Koopman Operator Meets Model Predictive Control. \textit{Automatica}, 93:149--160, 2018. DOI: \href{https://doi.org/10.1016/j.automatica.2018.03.046}{10.1016/j.automatica.2018.03.046}.
\bibitem{korda2018convergence} M. Korda and I. Mezić. On Convergence of Extended Dynamic Mode Decomposition to the Koopman Operator. \textit{Journal of Nonlinear Science}, 28:687--710, 2018. DOI: \href{https://doi.org/10.1007/s00332-017-9423-0}{10.1007/s00332-017-9423-0}.
\end{thebibliography}
\end{document}
